% Lösungen Einheit 3 von 3 – Sonderansätze und Modellkritik (zusammenbau v0.1)
\einheitenkopf{Einheit 3 von 3 · Sonderansätze und Modellkritik}

\erg{14}{a) $6$ – vom Hoch- zum Tiefpunkt ist es eine halbe Periode, zum nächsten Hochpunkt eine ganze. \quad b) $a = \frac{3 - (-3)}{2} = 3$; halbe Periode $3 - 1 = 2$, also $p = 4$ und $b = \frac{2\pi}{4} = \frac{\pi}{2} \approx 1{,}57$ \quad c) Viertelperiode $3$: $p = 12$ und $b = \frac{2\pi}{12} = \frac{\pi}{6} \approx 0{,}52$; $a = g(3) = f(3) = 2$ \quad d) halbe Periode $4$: $b = \frac{\pi}{4} \approx 0{,}79$; Fläche einer Halbwelle $\left[-\frac{a}{b} \mathrm{cos}(b x)\right]_0^{4} = \frac{2a}{b} = 8$, also $a = 4b \approx 3{,}14$ \quad e) Keine Durchschnittsänderung heißt $p(0) = p(40)$, der Scheitel liegt in der Mitte bei $x = 20$; Maximum $80 + 12 = 92$; $p(x) = a(x - 20)^2 + 92$ mit $p(0) = 400a + 92 = 80$, also $a = -0{,}03$ und $p(x) = -0{,}03(x - 20)^2 + 92$ \quad f) Waagerecht auf der $y$-Achse heißt $p'(0) = b = 0$, darum fehlt das lineare Glied. $c = 1{,}8$; eine Flächeneinheit sind $25$ m², die Fläche also $3{,}6$ Flächeneinheiten; mit der Nullstelle $x_0$ und $a = -\frac{c}{x_0^2}$ ist die Fläche $\frac{2}{3} c \cdot x_0 = 3{,}6$, also $x_0 = 3$ und $a = -0{,}2$ \quad g) Das Profil braucht an beiden Enden eine waagerechte Tangente; eine Parabel hat aber nur einen Scheitel, also nur eine Stelle mit waagerechter Tangente. Außerdem wechselt das Profil von einer Rechts- in eine Linkskrümmung, eine Parabel wechselt ihre Krümmung nie. Wer nur zeigt, dass eine Parabel nicht durch beide Punkte passt, hat nichts begründet.}
\erg{15}{a) $v = 100$ (halbe Länge); $s \cdot \left(e^{100/s} - e^{-100/s}\right) = 210$ mit dem Rechner: $s \approx 183{,}9$; $h(100) = 40$: $t = 40 - \frac{s}{2}\left(e^{100/s} + e^{-100/s}\right) \approx -171{,}8$}
\erg{16}{a) $p(0) = 3$, $p'(0) = -2$, $p(1) = 0{,}5$, $p'(1) = 1$; aus den ersten drei: $c = 3$, $b = -2$, $a = -0{,}5$; Probe: $p'(1) = 2a + b = -3$ – das ist nicht $1$, es gibt keine solche Funktion.}
\erg{17}{a) $q'(x) = -c \cdot x \cdot e^{-x^2/2}$: An den Rändern ist $x$ nicht null, $c$ ist positiv und die e-Potenz ist immer positiv – das Produkt ist dort nie null, für jede Wahl von $a$ und $c$. Ein einzelner Wert für $c$ zeigt das nicht.}
\erg{18}{a) Der Abstand von Hoch- und Tiefpunkt ist nur eine halbe Periode. Richtig: $p = 2 \cdot 4 = 8$ und $b = \frac{2\pi}{8} \approx 0{,}79$; $a = 3$. \quad b) Nach einer Periode wiederholt sich der Graph, der nächste Hochpunkt liegt also eine Periode weiter; der Tiefpunkt liegt wegen der Symmetrie des Sinusbogens genau in der Mitte dazwischen.}
